Nguồn: \href{https://vietnamnet.vn/nguyen-duy-phan-dinh-dieu-nhin-xa-de-hieu-gan-i5002796.html}{Báo VietNamNet}

\begin{quote}
Lời tòa soạn: "Nên là anh ấy nói đúng đấy, và thật là sâu sắc: Đúng là phải lùi xa mà nhìn lại, thì nó mới thật được, mới có thể bao quát, bao trùm được…" - Tròn 25 năm Internet vào Việt Nam, nhà thơ Nguyễn Duy nói về GS Phan Đình Diệu (1936-2018) -  nhà toán học, nhà khoa học máy tính nổi tiếng; Viện trưởng đầu tiên của Viện Khoa học Tính Toán \& Điều Khiển (nay là Viện Công nghệ Thông tin Việt Nam), đồng thời là Chủ tịch sáng lập Hội Tin học Việt Nam. GS Phan Đình Diệu được coi là người mở đường, "người anh cả" của ngành CNTT ở Việt Nam, cũng là "bức chân dung song hành" cùng nhà thơ Nguyễn Duy, thời Đổi mới.
\end{quote}


Ở vào giai đoạn “bản lề” của những năm Đổi mới, không hẹn mà cùng, họ từng hiện lên như hai bức chân dung song hành về tiếng nói phản biện bộc trực, đau đáu và chân thành trước thế cuộc, trước sự trì trệ bảo thủ và những khát vọng thay máu tư duy, cách nghĩ, cách làm… GS Phan Đình Diệu – nhà toán học, nhà khoa học máy tính nổi tiếng và nhà thơ Nguyễn Duy – tác giả của những bài thơ trữ tình thế sự từng gây chấn động một thời (Đánh thức tiềm lực, Kim Mộc Thuỷ Hoả Thổ, Nhìn từ xa Tổ quốc…).


Nhưng vẫn còn một điểm gặp thú vị nữa giữa họ: Trong thời chiến, nhà thơ Nguyễn Duy từng là một chiến sĩ thông tin; còn trong thời bình, GS Phan Đình Diệu lại chính là người đã thiết tha đề nghị Nhà nước mở đường cho Internet vào Việt Nam, tạo tiền đề quan trọng và có tính tiên quyết cho việc đưa ánh sáng của Internet tới Việt Nam, để điều kỳ diệu đó trở thành hiện thực.


\emph{Có phải sự tương đồng trong tiếng nói đã khiến hai ông gặp nhau, cả trên mặt chữ cũng như trong cuộc sống?}

Nhà thơ Nguyễn Duy: Thật ra khi mình được gặp anh Diệu mình cũng đã gần như… “xong việc” của mình rồi.

Như mình từng nói trong một bài phỏng vấn, rằng cánh cửa đổi mới vừa mới mở ra được một chút đã đóng sập lại, có nhiều người bị kẹt tay, mình cũng bị kẹt. Nhưng những gì cần nói ra, mình cũng đã kịp nói ra bằng hết. Chữ nghĩa dù có long đong lận đận, bằng cách này hay cách khác cũng đã lan được tới điểm này điểm nọ. Lúc ấy là ngồi chờ nó ngấm dần thôi.


Tương tự, khi mình và anh Diệu gặp nhau, cái phần anh ấy làm, anh ấy cũng đã làm rồi. Cái sự quý nhau lúc ấy nó là sự tâm đầu ý hợp, chứ không phải để bàn với nhau cùng làm cái này cái kia, không có chuyện đấy. Tư duy độc lập và gặp nhau thôi. 

\emph{Tôi nhớ là vào cái thời thơ Nguyễn Duy với những bài thơ dài bát ngát gây chấn động văn đàn và gần như không thể in cho tận mãi sau này, người ta thậm chí đã thu cả thơ ông vào băng cassette để chuyền tay nhau. Tương tự, vào cái thời internet chưa kịp về VN thì những bài phát biểu dậy sóng của GS Phan Đình Diệu cũng không dễ gì đến được tới rộng rãi quần chúng. Ông có cơ hội tiếp cận nhiều không?}

Ở thời điểm ấy mình quả thật không thể theo dõi tất cả những phát biểu, bài viết của anh Phan Đình Diệu, nhưng cơ bản là mình đã đọc được những bài quan trọng, cũng nhiều người đọc được, chuyền tay nhau đọc được. Đấy là một tiếng nói phản biện rất trung thực, rất ngay thẳng mà không hề ngoa ngôn, không hề dùng xảo ngữ; dũng cảm mạnh mẽ nhưng không sa vào chửi đổng, không phải để nói cho sướng mồm theo kiểu ngồi lê đôi mách (trong “Đánh thức tiềm lực”, mình cũng từng mỉa cái “nghề chửi đổng, nghề ngồi lê, nghề vu cáo…”). Tiếng nói ấy là thực tình, thực tâm muốn đóng góp cho đất nước, cho sự phát triển.

\emph{Trước và sau khi gặp GS Phan Đình Diệu, cảm nghĩ của ông về “bức chân dung song hành” có khác nhiều?}

Bên ngoài đời, anh Diệu rất là dịu dàng nhỏ nhẹ, anh hiền lắm, đôi khi ngây thơ ngơ ngác, mình nhớ anh hay hỏi: “Có thế thôi à?”. Nó gần như là ngược lại hoàn toàn với cái sắc sảo đanh thép trong chính luận của anh. Nhìn vẻ bề ngoài ấy, người ta không nghĩ là anh lại có tiếng nói mạnh mẽ, rắn rỏi đến thế. Từng từ ngữ mà anh Diệu sử dụng cho bài viết, bài nói của anh ấy rất là cẩn trọng, chính xác. Chính xác của ngôn ngữ, thuyết phục về phương pháp (là cả một khối kiến thức liên ngành cả về toán học, triết học, kinh tế học, xã hội học…), và trên hết là thái độ sống thẫm đẫm lòng trung thực.

\emph{“Có thế thôi à?” – Ngơ ngác đấy mà cũng là thất vọng đấy, là muốn khác, phải khác! Ông có nghĩ, cội nguồn của thẳng thắn phải chăng chính là sự hồn nhiên?}

Không, anh ấy hồn nhiên như toán học vậy thôi. Rất tự nhiên thoải mái.

\emph{Vậy ông nghĩ, cội nguồn đáng kể nhất ở đây là gì?}


Là sự chân chính. Anh Diệu rõ ràng là một người chân chính. Một người thực sự tốt, thực sự chân chính thì đứng ở đâu, làm gì, cũng là đều từ cái tốt ấy mà ra cả.

Bản chất anh Diệu là người tốt, lõi của anh là người tốt, lại còn là người tốt có tài, có tầm. Và vì mình đã xác định anh là người tốt nên mình tin tất cả những chuyện anh ấy làm bao giờ cũng xuất phát từ lòng tốt.

\emph{Một khi những tiếng nói thẳng thắn tìm đến nhau, thì phải chăng chính như ý thơ ông từng viết: “Những người tốt cần liên hiệp lại”?}

Đúng là những tiếng nói thẳng thắn luôn cần liên hiệp lại. Có thể có người nói trước, có người nói sau; lớp trước chưa làm được thì lớp sau làm tiếp; cái gì chưa ngấm ngay được về sau sẽ ngấm, từng chút từng chút một. Anh Diệu, mình nghĩ anh ấy có cái nhẫn nại đó, lòng tin đó. Thẳng thắn, và phải nói là vô cùng nhẫn nại.

{\bf “Bọn mình là những tế bào nhạy cảm của xã hội".}

\emph{Không ít những bóc tách, mổ xẻ trong những bài viết, bài phát biểu của GS Phan Đình Diệu, hay trong những bài thơ đau đáu nỗi nhân tình thế thái của Nguyễn Duy quả đã đưa tới những dự báo, cảnh báo từ rất sớm về những vấn nạn, mặt trái của sự trì trệ, thói đạo đức giả, tệ háo danh, tham nhũng… Đành rằng về sau, và nhất là những năm gần đây đã có những thay đổi tích cực, mạnh mẽ bằng vào những tác động “mưa dầm thấm lâu” hay những động thái quyết liệt của người cầm cương, nhưng ông có tiếc, giá như những tiếng nói đó được lắng nghe đúng thời điểm hơn?}

Quá tiếc ấy chứ! Nhưng biết làm thế nào được. Mình cũng phải chấp nhận cái sự ngấm dần, ngấm dần từng chút, được đến đâu hay đến đó, chứ không thể ngày một ngày hai mà chữa dứt điểm những căn bệnh nan y được.

\emph{Bên cạnh những rào cản bài xích, bảo thủ, những tiếng nói thẳng thắn vốn chưa bao giờ dễ nghe của Nguyễn Duy hay Phan Đình Diệu cũng đã gặp được sự tôn trọng, ủng hộ của những nhà lãnh đạo giàu chính kiến và biết tôn trọng chính kiến. Với ông, đó có là một sự khích lệ lớn? } 

Đối với mình mà nói, việc đầu tiên là cần nhìn nhận họ - những nhà lãnh đạo ấy ở góc độ một con người. Trong thời nào và ở giai tầng nào, cũng có những người tốt. Tốt cái đã! Những người tốt, ngay cả khi là thiểu số, nó cũng mang tới cho mình một niềm tin. Rằng, đến một lúc nào đó, phải, “những người tốt cần liên hiệp lại”…

 Không phải ai, ở đâu, lúc nào, trên cương vị nào cũng có thể tiện nói ra những điều họ thực nghĩ trong đầu, những điều họ thực lòng đồng cảm với mình. Nhưng quan trọng nhất, là họ biết mở lòng, lắng nghe những tiếng nói chân thành, trung thực. Bài “Đánh thức tiềm lực”, như các bạn đã biết, là “tiễn anh Sáu Dân đi làm kinh tế”, tức Thủ tướng Võ Văn Kiệt sau này. 

Mình dù có là nhà thơ, cái lõi của mình vẫn là thằng lính, là thằng nông dân, là một người dân từng trải qua chiến tranh và nghĩ về đất nước nhiều lắm, lung lắm! Sống đằm mình trong xã hội và chính bọn mình là tế bào nhạy cảm của xã hội, nên trong cái quan hệ giữa văn nghệ sỹ bọn mình và ông Võ Văn Kiệt hồi đó, nó hay lắm. Thời đó có thể nói ông Kiệt là một nhà lãnh đạo thực sự có một nhãn quan chân tình và sâu sắc, nên tất cả những phát biểu, những điều bọn mình nói với ông Sáu Dân, nó cứ thẳng băng. Qua đó, ông mới thấy được hết thực tế của đời sống xã hội. Chính bọn mình mới là người đưa lại cho ông thông tin chính xác về đời sống xã hội, chứ không phải là những báo cáo.

\emph{Khi là nhà khoa học Việt hiếm hoi được mời sang Mỹ trước cả khi Mỹ xoá bỏ lệnh cấm vận với Việt Nam, từ bờ tây Đại Tây Dương, GS Phan Đình Diệu từng viết câu thơ trải lòng mình với đất nước: “Bởi tự rất xa nhìn cái gần mới thật”. Một điểm gặp thú vị với tác giả “Nhìn từ xa Tổ quốc” - được viết lúc Nguyễn Duy sang Nga?}  

Về “nhìn xa” mà nói, thật ra anh Diệu có điều kiện tốt hơn bọn mình nhiều. Anh ấy là nhà khoa học nổi tiếng trong nước và cả quốc tế nữa, anh có một vị thế khác hẳn mình, một tâm thế khác hẳn mình; anh có điều kiện đi xa hơn cả về nghĩa đen lẫn nghĩa bóng. Anh được tiếp xúc sớm với những trí tuệ của nhân loại, lại còn là tiếp xúc trực tiếp (bọn mình dù sao thì cũng vẫn là tiếp xúc gián tiếp). Chính bởi thế mà cái nhìn của anh, tiếng nói của anh mới giàu tính dự báo và có tầm đến vậy.


Nên là anh ấy nói đúng đấy, và thật là sâu sắc: Đúng là phải lùi xa mà nhìn lại, nó mới thật được, mới có thể bao quát, bao trùm được. Độ lùi ấy có thể là không gian, cũng có thể là thời gian. Cũng như ngày hôm nay nhìn lại ngày hôm qua, nhìn lại những ấu trĩ ngột ngạt của một thời, rõ ràng mình thấy hôm nay tốt hơn cái thời mình làm những bài thơ ấy nhiều chứ! Hoàn toàn có cơ sở cho mình tin rằng đất nước này sẽ tốt hơn…

{\bf "Từng chút từng chút một, nhưng có thay đổi đấy!”/}

\emph{Vẫn còn một điểm gặp thú vị nữa giữa hai “bức chân dung song hành”: Trong thời chiến, nhà thơ Nguyễn Duy từng là một chiến sĩ thông tin; còn trong thời bình, GS Phan Đình Diệu lại chính là “người anh cả” của ngành CNTT ở Việt Nam. Là người từng thấm thía giá trị của thông tin trong thời chiến, ông nghĩ sao về những nỗ lực lớn của một người đã thiết tha đề nghị Nhà nước mở đường cho Internet vào Việt Nam giữa thời bình?}

Gọi là chiến sĩ thông tin cho “sang” nhưng mình thật ra chỉ là cái anh leo cột, nối dây ấy mà, gọi nôm na là lính đường dây cũng được. Sau này mình cũng có về làm Bộ Tư lệnh thông tin, cũng được tiếp xúc ít nhiều với kiến thức khoa học. Nhưng thông tin thời ấy của bọn mình nó đơn giản lắm, chủ yếu ở dạng truyền tin nội bộ thôi, chứ không nhiều chiều nhiều lớp như sau này. 


Internet nó lại là một vấn đề khác. Internet ở đây là mối quan hệ với thế giới. Vấn đề kỹ thuật Internet khác với tư tưởng Internet. Cái tin học, ngoài vấn đề khoa học ra thì vấn đề xã hội rất lớn, đấy là kết nối con người, kết nối các xã hội lại với nhau. Có Internet, mình mới có được sự cởi mở như ngày hôm nay, nếu không thì mình tắc tử rồi. Công đấy của anh Diệu, mới là!

\emph{Ông có tiếc, kể mà Internet đến sớm hơn, thì thơ Nguyễn Duy đã được đọc nhiều hơn, kịp thời hơn, và tiếng nói Phan Đình Diệu, cũng đã dễ dàng phủ sóng hơn, chỉ cần bằng một cú nhấp chuột? Hay thật ra, người cần lắng nghe nhất ở đây, lại là thiểu số, hơn là đa số?}

Nhưng cái thiểu số đó, nó có sức lan tỏa của nó. Nó đi đến tất cả, nó đi đến từng điểm. Mỗi điểm như thế, nó đều có công chúng, có từ trường của của nó, có không gian lan tỏa của nó. Những điều mình viết, những điều anh Diệu nói hay viết, không phải mọi người đều hiểu, không phải đều tiếp nhận. Nó tới được từng điểm thôi. Điểm quan trọng là những người có trách nhiệm, có tri thức, có trí tuệ. Từ đó, nó sẽ mở dần, lan tỏa dần. Có thể rất chậm nhưng nó có những vòng tròn đồng tâm của nó.

Ở đây mình có hai cơ sở để mà tin. Thứ nhất là cái sự thay đổi trong lòng người, nó có thể rất chậm, từng chút từng chút một, không dễ gì một sớm một chiều mà thậm chí phải qua nhiều thế hệ. Nên không thể sốt sắng mà được, phải kiên trì. Hai là, ở thời hội nhập này, không thể đơn độc mà đi được, mà rõ ràng phải đi theo cái mạch chung của thế giới, của nhân loại. Cái bên trong phải phù hợp với cái bên ngoài, phải có một đường dẫn, đường thông ra tới bên ngoài. Hai cái đó, nó sẽ làm cho đất nước này thay đổi. Và cho đến bây giờ, mặc dù sự trì trệ vẫn còn đầy ra đấy, nhưng mà cũng phải thừa nhận, xã hội đã từng chút từng chút một thay đổi theo chiều hướng tích cực hơn, gần với quy luật phát triển chung của nhân loại hơn.

\emph{Đã bao giờ ông bị giằng xé giữa khát vọng được nói ra bằng hết, và cố gắng kìm nén lại để đỡ làm liên luỵ đến những người “đồng thanh tương ứng” với mình, và hơn hết là vì chính sự an toàn của mình? Đã bao giờ giữa những cuộc chuyện, mà hai ông từng nhắc tới sự giằng xé đó?}

Không, chẳng bao giờ. Nhưng bọn mình cùng ngầm hiểu, cái sợi dây gắn kết giữa bọn mình, cái làm nên sự tri âm tri kỷ ở đây, chính là cái sự “liều mình như chẳng có” ấy, mà cái gốc của nó, là đều từ tấm lòng đối với xã hội, với đất nước, với sự phát triển chung… Anh Diệu nếu như chỉ cắm cúi làm khoa học thuần túy mà không để tâm gì đến những vấn đề của thời cuộc thì anh ấy đã ở vào một cương vị khác, cũng không ai trách anh ấy cả. Còn một khi anh tình nguyện mang vác cái trách nhiệm đó, thì có nghĩa sẽ phải chấp nhận va chạm và đằng sau sự va chạm đấy có thể là sự mất mát, thua thiệt.

Khi đưa tiếng nói đó vào thơ, mình cũng nghĩ, đóng góp được gì thì đóng góp thôi, làm được cái gì có ích, có ý nghĩa thì mình làm thôi, còn tác dụng đến đâu thì mình không tính được. Mình tin là những người hiểu biết, người ta cũng sẽ đồng cảm với mình, từng chút, từng chút một người nọ sẽ lan truyền cho người kia, mình tin là nó có tác dụng chứ không phải không đâu. Nó dĩ nhiên không có tác dụng trực tiếp vào việc giúp thay đổi thể chế này nhưng mà nó có tác dụng trong sự thay đổi về tâm tư của con người, của xã hội. Chính từ đó, nó sẽ tác động đến ý thức.

Thật sự là mình hài lòng với những gì mình đã làm và cũng không nghĩ lúc đó mình làm được thế. Khi ở trong hoàn cảnh ấy thật ra không có đủ thời gian, điều kiện để tự quan sát mình đâu, thấy việc thì làm, thích thì làm, còn việc nó tới đâu, nó có thành cái gì không, hay nó có hại gì không thì chỉ cảm thấy bâng quơ thế thôi, chứ không để bị phụ thuộc vào nó. Mình tin anh Diệu hẳn cũng nghĩ thế thôi. Khi mình chân thành, thẳng thắn, mình sẽ luôn nghĩ tới cái đích đến tốt đẹp của nó. Nó chắc chắn sẽ tốt dần lên, tuy chậm nhưng vẫn có tiến triển, rồi đến lúc nào đó, cái tốt sẽ nhiều hơn, được phổ biến hơn. Cũng không bao giờ nó tốt hoàn toàn cả…

Dẫu có sao thì cũng đừng có thở dài…

\emph{Khi những bài thơ, bài viết bị cho là khó in, khó đăng trong điều kiện xuất bản trong nước, cả ông và GS Phan Đình Diệu đều chọn một cách ứng xử giống nhau: Không tìm kiếm cơ hội thay thế ở bên ngoài biên giới hình chữ S. Với ông là vì sao?}

Thường khi viết ra được một bài thơ mà mình tâm huyết, mình luôn muốn được công bố ngay, muốn nó đi vào được đám đông ngay thay vì nghĩ đến hậu quả nào sẽ xảy tới với mình. Nhưng đó là công chúng nào, lại cần thận trọng. Hồi mình qua Mỹ, từng có một tạp chí hải ngoại hỏi in một bài thơ không được đăng trong nước của mình lúc ấy, nhưng mình từ chối. Vì nếu mình đưa cho họ in, lại càng khó để được in trong nước, trong khi cái mình muốn là thơ của tôi, tiếng nói của tôi phải đến được với dân tôi, nước tôi, chứ nếu chỉ để nói đổng bên ngoài thì uổng lắm. Chưa có bài nào mình cho nước ngoài in trước cả, bao giờ cũng phải đợi trong nước in đã. Trong chuyện này, mình tin anh Diệu cũng nghĩ giống mình.

\emph{Nói thẳng thì thường khó nghe, đi trước thời cuộc thì thường cô độc. Ông còn cảm giác cô độc đó không, khi nhìn vào “bức chân dung song hành”?}

Cũng phải nhờ sự cô đơn, cô độc ấy thì mới cất lên tiếng nói được. Những ý nghĩ ấy thật ra rất nhiều người nghĩ mà người dám phát ngôn lại rất ít. Ngồi thì thào với nhau lúc trà dư tửu hậu rất là sôi nổi, nhưng ra trước công luận, hầu như mọi người lại im lặng, nén đi.

Cho nên, mình quý anh Diệu nhất là ở cái tiếng nói đó, ở cái cách mà anh lên tiếng. Một tiếng nói trung thực của một người có tài, có tâm, có tầm, có bản lĩnh, có nhân cách và tri thức.

Mình trọng anh Diệu là trọng ở cái thực chất mà anh ấy đóng góp cho xã hội, cho đất nước, chứ không phải là danh tiếng giáo sư của anh. Trong “Kim Mộc Thuỷ Hoả Thổ” mình thậm chí còn từng “bái phục giáo sư vặt lông vịt thiên tài” kia mà! Có một thời, cái thành phần “giáo sư vặt lông vịt” rất là nhiều, rất phổ biến, nhưng mà một Giáo sư như là GS Phan Đình Diệu thì thật sự là hiếm hoi, thật sự là đáng kính…

Lê Quân (thực hiện)